\chapter{GraphicsDevice and the Rendering Pipeline}
\label{ch:graphicsdevice}

\cnaclass{GraphicsDevice} is the hub every other class in this Part draws through. This
chapter covers its own 887-line header directly, then the resource base class
(\cnaclass{GraphicsResource}) every texture, buffer, and state object inherits, and the
disposal ordering discipline that keeps a multi-backend rendering pipeline from ever calling
into an already-destroyed graphics context.

\section{Construction: one real constructor, one CNA addition}

\cnaclass{GraphicsDevice} inherits \texttt{System::Object} and \texttt{System::IDisposable},
is non-copyable and non-movable, and offers exactly two constructors: the real XNA one,
\texttt{GraphicsDevice(GraphicsAdapter\&, GraphicsProfile, const PresentationParameters\&)},
and a \texttt{NOXNA} parameterless \texttt{GraphicsDevice()} for headless/no-window contexts.
Six events match FNA's surface exactly: \texttt{Disposing}, \texttt{DeviceLost} (never
actually raised on desktop), \texttt{DeviceReset}, \texttt{DeviceResetting},
\texttt{ResourceCreated}, \texttt{ResourceDestroyed}.

\section{The property surface}

Nineteen properties cover the device's full observable state: \texttt{IsDisposed},
\texttt{GraphicsDeviceStatus}, \texttt{Adapter}, \texttt{GraphicsProfile},
\texttt{PresentationParameters}, \texttt{DisplayMode}, \texttt{Textures},
\texttt{SamplerStates}, \texttt{VertexTextures}, \texttt{VertexSamplerStates},
\texttt{BlendState}, \texttt{DepthStencilState}, \texttt{RasterizerState},
\texttt{ScissorRectangle}, \texttt{Viewport}, \texttt{BlendFactor}, \texttt{MultiSampleMask},
\texttt{ReferenceStencil}, \texttt{Indices}. One internal detail worth flagging up front,
because it recurs throughout Part~IV: \texttt{deviceStatus\_} is only ever driven away from
\texttt{Normal} by the \texttt{D3D9} backend's real device-lifecycle callback
(Chapter~\ref{ch:direct3d-backends}) --- every other backend never touches it, so
\texttt{GraphicsDeviceStatus} stays \texttt{Normal} for the whole life of the process on any
other backend.

\section{Clear, Present, Reset}

\texttt{Clear(Color)} and the full \texttt{Clear(ClearOptions, Color, float depth, int
stencil)} are both present, alongside two unmarked (arguably under-tagged) CNA convenience
overloads, \texttt{Clear(float,float,float,float)} and \texttt{Clear(Color, float depth)}.
FNA's \texttt{Clear(ClearOptions, Vector4, float, int)} overload (color as a
\texttt{Vector4} rather than a \texttt{Color}) is the one XNA overload genuinely missing.
\texttt{Present()} exists only in its no-argument form --- FNA's 3-argument
\texttt{Present(Rectangle?, Rectangle?, IntPtr)} overload (source/destination rectangles and
an override window handle) is absent. \texttt{Reset()} comes in three forms (no args, new
\cnaclass{PresentationParameters}, new parameters plus a new \cnaclass{GraphicsAdapter}) plus
a fourth NOXNA pointer-adapter overload.

\begin{sourcenote}
\textbf{From the FNA-fidelity audit (\texttt{docs/graphicsdevice-fna-audit.md}):} the three
concrete gaps against FNA's public surface are the \texttt{Present} 3-arg overload, the
\texttt{Clear(ClearOptions, Vector4, ...)} overload, and \texttt{GetRenderTargetsNoAllocEXT}.
Several convenience overloads and aliases (the two extra \texttt{Clear} forms above, plus
\texttt{SetIndexBuffer}/\texttt{GetIndexBuffer}/\texttt{GetVertexBuffer} aliases) exist without
their expected \texttt{NOXNA} tag --- a real, specific instance of
Chapter~\ref{ch:design-philosophy}'s \texttt{NOXNA}-discipline being imperfectly applied, worth
knowing about if you are relying on the \texttt{NOXNA} build check to catch every non-XNA
member mechanically.
\end{sourcenote}

\section{Draw calls: concrete overloads where FNA has generics}

\texttt{DrawPrimitives}, \texttt{DrawIndexedPrimitives}, and \texttt{DrawInstancedPrimitives}
are all present with FNA-matching signatures. FNA's \texttt{DrawUserPrimitives<T>}/%
\texttt{DrawUserIndexedPrimitives<T>} generic methods --- which in C\# work against any
\texttt{struct} implementing \texttt{IVertexType} --- become, in CNA, a fixed set of
concretely-typed overloads (\texttt{VertexPositionColor}, \texttt{VertexPositionColorTexture},
\texttt{VertexPositionTexture}, \texttt{VertexPositionNormalTexture}), plus a raw
\texttt{void*} fallback and an explicit-\cnaclass{VertexDeclaration} overload for anything
outside that fixed set. This is a direct consequence of C\texttt{++} template instantiation
requiring concrete call sites rather than C\#'s runtime generic dispatch --- and is the same
underlying constraint behind \texttt{GetBackBufferData}'s narrowing to \texttt{Color} only
(no generic pixel-format element type).

\section{Resource tracking and disposal ordering}
\label{sec:resource-lifetime}

\cnaclass{GraphicsDevice} maintains a private, non-owning
\texttt{std::vector<GraphicsResource*> resources\_}, mutated only through
\texttt{AddResourceReference}/\texttt{RemoveResourceReference} --- both NOXNA, both existing
purely so the device can guarantee a specific disposal order. \texttt{ResourceCreated} fires
immediately after registration; \texttt{ResourceDestroyed} fires \emph{before}
de-registration, carrying a \texttt{Name}/\texttt{Tag} snapshot captured while the resource is
still valid. Removal is swap-and-pop (O(1), order not preserved) --- deliberate, since nothing
about resource tracking depends on iteration order.

The disposal sequence this buys is the single most important invariant in this chapter:
calling \texttt{GraphicsDevice::Dispose()} copies and clears \texttt{resources\_}, disposes
every tracked resource \emph{first}, and only then tears down the native backend context
(SDL/Vulkan/GL). No graphics-API call can ever be issued against an already-destroyed
context, because every resource that might issue one is disposed strictly before the context
that would receive it.

\begin{sourcenote}
\textbf{A documentation trap worth naming explicitly.} \texttt{docs/graphicsresource-fna-audit.md}
states, as an open gap, that ``device resource tracking: not implemented'' and that a
\texttt{GraphicsDeviceResetting()} callback is missing. Both claims are stale --- the resource
tracking mechanism described in this section, and the \texttt{DeviceResetting} event on
\cnaclass{GraphicsDevice} itself, both demonstrably exist in the current header. This looks
like exactly the kind of documentation drift Chapter~\ref{ch:what-is-cna} warned about:
treat a specific, checkable claim in \texttt{docs/} with the same ``verify against current
source, not memory'' discipline the project asks of itself, even when the claim comes from
the project's own documentation.
\end{sourcenote}

\subsection{Per-backend disposal caveats}

The ordering guarantee above is a CNA-wide invariant, but what actually happens to a native
GPU handle on disposal is backend-specific, and \texttt{docs/graphics-resource-lifetime.md}
is candid about the sharp edges: on \textbf{EasyGL}, disposing a bound
\cnaclass{RenderTarget2D} without first unbinding it can leave a dangling framebuffer
attachment, and there is currently no GL context-loss recovery; on \textbf{Vulkan},
destroying a render target still referenced by an in-flight frame is undefined behavior ---
there is no \texttt{vkDeviceWaitIdle} guard at that call site; on \textbf{Bgfx},
\texttt{bgfx::destroy()} only \emph{queues} destruction for the next \texttt{bgfx::frame()},
so \texttt{bgfx::shutdown()} must never run before every resource backend has actually been
destroyed. None of these are exercised by ordinary game code following the
\texttt{Begin}/\texttt{Draw}/\texttt{End} pattern from earlier chapters, but they matter if
your own code manages render-target lifetime manually across frame boundaries.

\section{GraphicsResource: the shared base}

Every texture, buffer, and state object in this Part inherits \cnaclass{GraphicsResource},
which itself inherits \texttt{System::Object}/\texttt{System::IDisposable} and provides
\texttt{Disposing}, a get-only \texttt{GraphicsDevice}, \texttt{IsDisposed}, a virtual
\texttt{Name}, and a raw, caller-owned \texttt{Tag} (\texttt{System::Object*}).
\texttt{ToString()} returns \texttt{Name} if set, otherwise the type name. Its protected
constructor, \texttt{GraphicsResource(GraphicsDevice* device = nullptr)}, accepts a null
device deliberately --- the four state-object classes in Chapter~\ref{ch:state-objects} each
ship static preset instances (\cnaclass{BlendState::Opaque},
\cnaclass{RasterizerState::CullNone}, and so on) that exist before any real device does, and
those presets simply skip tracking/registration/events entirely.

A subtlety worth internalizing: GPU handles are freed on \texttt{Dispose()}, not on the
C\texttt{++} destructor. The destructor does call \texttt{Dispose(false)} as a safety net,
but that path deliberately skips the \texttt{Disposing}/\texttt{ResourceDestroyed} events ---
matching the standard C\# \texttt{IDisposable} finalizer-vs-explicit-dispose distinction, now
expressed in C\texttt{++} terms. Move semantics transfer \texttt{unique\_ptr}-owned GPU state
correctly but do \emph{not} update the device's tracking-list pointer entry, which means
moving a resource that is currently tracked by a live \cnaclass{GraphicsDevice} to a new
memory address is unsafe and unsupported --- construct in place, or move before the resource
is ever registered.

\section{GraphicsAdapter, DisplayMode, and the exception types}

\cnaclass{GraphicsAdapter} exposes \texttt{CurrentDisplayMode}, \texttt{SupportedDisplayModes},
\texttt{Description}, and a static, re-evaluated-every-call \texttt{DefaultAdapter} --- do not
cache it across a call to the static \texttt{AdaptersChanged()}. Two properties are honestly
limited on this platform: \texttt{DeviceId}/\texttt{VendorId} are queried through Linux
\texttt{sysfs} where available and return \texttt{0} otherwise, while
\texttt{Revision}/\texttt{SubSystemId} are always \texttt{0} --- SDL3 has no query surface for
either. \cnaclass{DisplayMode} pairs \texttt{Width}/\texttt{Height}/\texttt{Format} with a
computed \texttt{AspectRatio}; \cnaclass{DisplayModeCollection} supports both integer indexing
and a \cnaclass{SurfaceFormat}-filtered \texttt{operator[]}. \cnaclass{GraphicsDeviceStatus}
is the three-value enum (\texttt{Normal}, \texttt{Lost}, \texttt{NotReset}) already introduced
above; \cnaclass{GraphicsProfile} is the familiar XNA \texttt{Reach}/\texttt{HiDef} pair, whose
\emph{only} backend that actually consults a real device capability structure for it is
\texttt{D3D9} (Chapter~\ref{ch:direct3d-backends}) --- everywhere else the profile is present
in the API but decorative. \cnaclass{DeviceLostException},
\cnaclass{DeviceNotResetException}, and \cnaclass{NoSuitableGraphicsDeviceException} are all
thin \texttt{std::runtime\_error} subclasses with default and custom-message constructors.

\section{ClearOptions: one confirmed cross-backend gap}

\cnaclass{ClearOptions} is a three-value flags enum (\texttt{Target=1}, \texttt{DepthBuffer=2},
\texttt{Stencil=4}) with the full bitwise operator set hand-implemented, the same pattern seen
on \cnaclass{DisplayOrientation} in Chapter~\ref{ch:game-loop}. One flag is a confirmed,
still-open gap at the time of writing: \texttt{ClearOptions::Stencil} is a no-op on EasyGL and
Vulkan (Task~871), with Bgfx unverified --- clearing the stencil buffer this way currently
does nothing on the two backends where it has actually been checked. This is a small, concrete
example of the kind of specific, checkable claim Chapter~\ref{ch:what-is-cna} promised this
book would preserve rather than round off to ``clearing works.''
